\subsection{01-Trie}

把数字按二进制从高位到低位插入的二叉树（每位 0/1 一个儿子，结点计数为其子树中元素个数），可完成与平衡树相同的多重集操作，且实现与常数更小。视为多重集：重复值可存，\texttt{remove(x)} 仅删除一个等于 $x$ 的元素（不存在则无事发生）。值域需为非负整数：固定 $H=30$ 支持 $0\le x\le 2^{31}-1$，更大值域改大 \texttt{H}。

各操作单次 $\mathcal{O}(\log V)$（$V$ 为值域）：\texttt{insert(x)} 插入一个 $x$；\texttt{remove(x)} 删除一个等于 $x$ 的元素；\texttt{rank(x)} 返回严格小于 $x$ 的元素个数 $+1$（即插入 $x$ 后 $x$ 得到的最小排名），$x$ 不存在时返回 $-1$；\texttt{index(k)} 返回第 $k$ 小元素（下标从 1 开始）；\texttt{contains(x)}、\texttt{size()} 分别判断存在性与元素个数。

\begin{lstlisting}
struct Trie01{
    using u32 = unsigned int;
    static constexpr int H = 30;
    std::vector<u32> cnt;
    std::vector<std::array<u32,2>> son;
    Trie01(){
        newNode();
    }
    void insert(int x){
        u32 u = 0;
        cnt[u]++;
        for (int b = H;b >= 0;b--) {
            u32 bit = (x >> b) & 1;
            if (son[u][bit] == 0) son[u][bit] = newNode();
            u = son[u][bit];
            cnt[u]++;
        }
    }
    void remove(int x){
        if (!contains(x)) return;
        u32 u = 0;
        cnt[u]--;
        for (int b = H;b >= 0;b--) {
            u = son[u][(x >> b) & 1];
            cnt[u]--;
        }
    }
    int rank(int x) const {
        u32 less = 0,u = 0;
        for (int b = H;b >= 0;b--) {
            u32 bit = (x >> b) & 1;
            if (bit && son[u][0]) less += cnt[son[u][0]];
            if (son[u][bit] == 0) return -1;
            u = son[u][bit];
        }
        if (cnt[u] == 0) return -1;
        return less + 1;
    }
    bool contains(int x) const {
        u32 u = 0;
        for (int b = H;b >= 0;b--) {
            u = son[u][(x >> b) & 1];
            if (u == 0) return 0;
        }
        return cnt[u] != 0;
    }
    int index(u32 rk) const {
        u32 u = 0,ans = 0;
        for (int b = H;b >= 0;b--) {
            u32 c0 = son[u][0] ? cnt[son[u][0]] : 0;
            if (rk <= c0) u = son[u][0];
            else {
                rk -= c0;
                ans |= 1u << b;
                u = son[u][1];
            }
        }
        return ans;
    }
    u32 size() const {
        return cnt[0];
    }
private:
    u32 newNode(){
        cnt.emplace_back(0);
        son.emplace_back(std::array<u32,2>({0,0}));
        return cnt.size() - 1;
    }
};
\end{lstlisting}
